% GENERATED FILE. Edit canonical lessons, then run npm run generate:books.
% canonical-source-hash: fnv1a64:155c6665fbbd6cef
% canonical-lessons: RU-C124-purple, RU-C124-brown, RU-C124-grey, RU-C124-orange, RU-C124-busy

\chapter{Purple, Brown, Grey, Orange, Busy}
\label{ch:ru-purple-brown-grey-orange-busy}
\hlchaptermodality{124}

\begin{chapteropening}
\textbf{By the end of this chapter:} \emph{I can say purple, brown, grey, orange and busy.}

Complete the last lesson of chapter 124: I can say purple, brown, grey, orange and busy.
\end{chapteropening}

\section[\texorpdfstring{\ru{фиолетовый} (fiolétovyy)}{фиолетовый (fiolétovyy)}]{\ru{фиолетовый} (fiolétovyy) — purple}

\label{lesson:RU-C124-purple}



\begin{quote}
Before the new one: say the Russian for stairs, then the Russian for pink.
\end{quote}

\subsection*{You'll want to know: \ru{фиолетовый}}
\textbf{\ru{фиолетовый}} (\emph{fiolétovyy}) — \textquotedblleft{}purple\textquotedblright{}.

\begin{cousinweb}
\textbf{\ru{фиолетовый}} — its \emph{\ru{ф}-} shows it came through German \emph{violett} from French \emph{violet}, from \emph{violette}, Latin \emph{viola}, \textquotedblleft{}the violet\textquotedblright{} (Vasmer).
\end{cousinweb}

\begin{grammarlens}[title={What you've built}]
One more word for this chapter.
\end{grammarlens}

\subsection*{Your turn}
\begin{itemize}
\raggedright
  \item \emph{Say it:} \emph{fiolétovyy}
  \item \emph{Say it:} \emph{fiolétovyy}, once more
  \item \emph{Say it:} say \emph{fiolétovyy}
  \item \emph{Recall:} say the Russian for stairs, then the Russian for pink, then say \emph{fiolétovyy} again
\end{itemize}

\begin{tcolorbox}[breakable,colback=teal!4,colframe=teal!35!black,title={Before you move on}]
What does \ru{фиолетовый} mean? (\textquotedblleft{}Purple\textquotedblright{}.) Say it once more.
\end{tcolorbox}
\section[\texorpdfstring{\ru{коричневый} (koríchnevyy)}{коричневый (koríchnevyy)}]{\ru{коричневый} (koríchnevyy) — brown}

\label{lesson:RU-C124-brown}



\begin{quote}
Before the new one: say the Russian for pink, then the Russian for purple.
\end{quote}

\subsection*{You'll want to know: \ru{коричневый}}
\textbf{\ru{коричневый}} (\emph{koríchnevyy}) — \textquotedblleft{}brown\textquotedblright{}.

\begin{cousinweb}
\textbf{\ru{коричневый}} is \textquotedblleft{}cinnamon-coloured\textquotedblright{}, from \emph{\ru{корица}}, \textquotedblleft{}cinnamon\textquotedblright{}, a diminutive of \emph{\ru{кора}}, \textquotedblleft{}bark\textquotedblright{} (Vasmer): cinnamon is a bark.
\end{cousinweb}

\begin{grammarlens}[title={What you've built}]
One more word for this chapter.
\end{grammarlens}

\subsection*{Your turn}
\begin{itemize}
\raggedright
  \item \emph{Say it:} \emph{koríchnevyy}
  \item \emph{Say it:} \emph{koríchnevyy}, once more
  \item \emph{Say it:} say \emph{koríchnevyy}
  \item \emph{Recall:} say the Russian for pink, then the Russian for purple, then say \emph{koríchnevyy} again
\end{itemize}

\begin{tcolorbox}[breakable,colback=teal!4,colframe=teal!35!black,title={Before you move on}]
What does \ru{коричневый} mean? (\textquotedblleft{}Brown\textquotedblright{}.) Say it once more.
\end{tcolorbox}
\section[\texorpdfstring{\ru{серый} (séryy)}{серый (séryy)}]{\ru{серый} (séryy) — grey}

\label{lesson:RU-C124-grey}



\begin{quote}
Before the new one: say the Russian for purple, then the Russian for brown.
\end{quote}

\subsection*{You'll want to know: \ru{серый}}
\textbf{\ru{серый}} (\emph{séryy}) — \textquotedblleft{}grey\textquotedblright{}.

\begin{cousinweb}
\textbf{\ru{серый}} $\leftarrow$ Proto-Slavic \emph{*sěrŭ}, from \emph{*xoiro-} — inherited. Vasmer relates it to Old Norse \emph{hárr}, \textquotedblleft{}grey, hoary\textquotedblright{}, and Old English \emph{hār}, so English \emph{hoar} is a cousin.
\end{cousinweb}

\begin{grammarlens}[title={What you've built}]
One more word for this chapter.
\end{grammarlens}

\subsection*{Your turn}
\begin{itemize}
\raggedright
  \item \emph{Say it:} \emph{séryy}
  \item \emph{Say it:} \emph{séryy}, once more
  \item \emph{Say it:} say \emph{séryy}
  \item \emph{Recall:} say the Russian for purple, then the Russian for brown, then say \emph{séryy} again
\end{itemize}

\begin{tcolorbox}[breakable,colback=teal!4,colframe=teal!35!black,title={Before you move on}]
What does \ru{серый} mean? (\textquotedblleft{}Grey\textquotedblright{}.) Say it once more.
\end{tcolorbox}
\section[\texorpdfstring{\ru{оранжевый} (oránzhevyy)}{оранжевый (oránzhevyy)}]{\ru{оранжевый} (oránzhevyy) — orange (the colour)}

\label{lesson:RU-C124-orange}



\begin{quote}
Before the new one: say the Russian for brown, then the Russian for grey.
\end{quote}

\subsection*{You'll want to know: \ru{оранжевый}}
\textbf{\ru{оранжевый}} (\emph{oránzhevyy}) — \textquotedblleft{}orange (the colour)\textquotedblright{}.

\begin{cousinweb}
\textbf{\ru{оранжевый}} is a late borrowing from French \emph{orange}, \textquotedblleft{}an orange\textquotedblright{}, perhaps through German; the source is Persian and Arabic \emph{nāranj}, whose home is sought in India, Sanskrit \emph{nāraṅga-} (Vasmer).
\end{cousinweb}

\begin{grammarlens}[title={What you've built}]
One more word for this chapter.
\end{grammarlens}

\subsection*{Your turn}
\begin{itemize}
\raggedright
  \item \emph{Say it:} \emph{oránzhevyy}
  \item \emph{Say it:} \emph{oránzhevyy}, once more
  \item \emph{Say it:} say \emph{oránzhevyy}
  \item \emph{Recall:} say the Russian for brown, then the Russian for grey, then say \emph{oránzhevyy} again
\end{itemize}

\begin{tcolorbox}[breakable,colback=teal!4,colframe=teal!35!black,title={Before you move on}]
What does \ru{оранжевый} mean? (\textquotedblleft{}Orange (the colour)\textquotedblright{}.) Say it once more.
\end{tcolorbox}
\section[\texorpdfstring{\ru{занятый} (zányatyy)}{занятый (zányatyy)}]{\ru{занятый} (zányatyy) — busy}

\label{lesson:RU-C124-busy}



\begin{quote}
Before the new one: say the Russian for grey, then the Russian for orange.
\end{quote}

\subsection*{You'll want to know: \ru{занятый}}
\textbf{\ru{занятый}} (\emph{zányatyy}) — \textquotedblleft{}busy\textquotedblright{}.

\begin{cousinweb}
\textbf{\ru{занятый}} is \textquotedblleft{}taken up\textquotedblright{}: \emph{\ru{за}-} with \emph{-\ru{нятый}}, from \emph{-\ru{нять}}, the old Slavic verb \emph{*jęti}, \textquotedblleft{}to take\textquotedblright{}, as in \emph{\ru{взять}}, \textquotedblleft{}to take\textquotedblright{}, and \emph{\ru{понимать}}, \textquotedblleft{}to understand\textquotedblright{} (Vasmer).
\end{cousinweb}

\begin{grammarlens}[title={What you've built}]
One more word for this chapter.
\end{grammarlens}

\subsection*{Your turn}
\begin{itemize}
\raggedright
  \item \emph{Say it:} \emph{zányatyy}
  \item \emph{Say it:} \emph{zányatyy}, once more
  \item \emph{Say it:} say \emph{zányatyy}
  \item \emph{Recall:} say the Russian for grey, then the Russian for orange, then say \emph{zányatyy} again
\end{itemize}

\begin{tcolorbox}[breakable,colback=teal!4,colframe=teal!35!black,title={Before you move on}]
What does \ru{занятый} mean? (\textquotedblleft{}Busy\textquotedblright{}.) Say it once more.
\end{tcolorbox}
